\section{Independence-Aware Reinforcement Learning}\label{sec:method}

\subsection{The independence certificate}\label{sec:certificate}
Let $G$ be the undirected bipartite graph on the places and transitions of the net,
with an edge between a place and every transition that consumes from it or produces
into it. Transitions that can never fire (their input places can never be marked)
are removed first; this least-fixpoint pruning over-approximates what can fire, so it
is sound. The connected components $\comp_1,\dots,\comp_K$ of the remaining graph
partition the places and transitions; every action, event and reward belongs to
exactly one component. Any shared place, whether a queue or a resource pool, puts its
transitions in the same component. The partition depends only on the net, not on the
marking or on the policy. For the insurer of Sect.~\ref{sec:setup}, each process of
each region is one component; if claims and complaints share their clerks, they form
one component.

\subsection{Factored credit: Net-Factored GAE}\label{sec:nfgae}
For a component $c$, let $k^c_1<k^c_2<\dots$ be the decision epochs at which an action
of $c$ is chosen, let $R^c_i$ be the reward of $c$'s events in
$[\tau_{k^c_i},\tau_{k^c_{i+1}})$ and $\Gamma^c_i = e^{-\beta(\tau_{k^c_{i+1}}-\tau_{k^c_i})}$.
\nfgae{} runs GAE per component, on the component's own decision clock:
\begin{equation}
  \delta^c_i = R^c_i + \Gamma^c_i V_c(s_{k^c_{i+1}}) - V_c(s_{k^c_i}), \qquad
  \hat A^c_i = \sum_{l\ge 0} \Big(\lambda^l \prod_{j<l}\Gamma^c_{i+j}\Big)\, \delta^c_{i+l},
\end{equation}
and each decision is updated with the advantage of its own component. The critic $V_c$
pools only the action nodes of the deciding component. With $K=1$, \nfgae{} is exactly
PPO with SMDP-GAE.

\subsection{Component-scoped postponement}\label{sec:postpone}
Instead of one global postpone pseudo-action, every component with an enabled action
gets its own. Choosing it blocks only that component: its actions are disabled until
one of \emph{its own} events fires, while the other components keep deciding. Time
advances to the next event once every component has acted, has nothing enabled, or
is waiting. A component's decision instants therefore depend only on its own state,
which preserves the independence the certificate relies on. Global postponement does
not: when one process should wait, all processes are frozen, so a process's decision
instants depend on the state of the others. \todo{counterexample figure or one-line
example}

\subsection{Component turns}\label{sec:turns}
When several components have enabled work at the same instant, they are offered one at
a time, lowest component index first, and the decision chooses among the active
component's bindings only. No simulated time passes between turns, components share no
place, and the setting is work-conserving or component-postponing, so the order of turns
does not change any component's outcomes. Each decision then observes and encodes only
its own component's subgraph, so the cost of a decision no longer grows with the size of
the organisation.

\subsection{Guarantees}\label{sec:guarantees}
We state the results informally; proofs are in the supplementary material.

\noindent\todo{formal statements; proofs in the supplement}
\begin{itemize}
  \item \textbf{Soundness.} If (C1) components are place-disjoint (by construction),
    (C2) the logit of a binding depends only on its own component's marking (true for
    the policy network of Sect.~\ref{sec:network}, whose message passing follows the
    arcs of the net), and (C3) the agent postpones per component or not at all, then
    the components evolve conditionally independently given the state.
  \item \textbf{Unbiasedness.} Under these conditions the \nfgae{} policy gradient is
    unbiased for the organisation's objective $J=\sum_c J_c$, and with $K=1$ it
    coincides with PPO.
  \item \textbf{Variance.} For any state-dependent baseline, \nfgae{} removes the term
    $\mathbb{E}\big[\lVert\nabla\log\pi\rVert^2\,\mathrm{Var}(\text{other components'
    return}\mid s)\big]$ from the variance of the policy gradient. No state-dependent
    baseline can remove this term, since it stems from the other components' future
    randomness.
  \item \textbf{Component turns} leave every component's dynamics unchanged and give the
    active component the same logits as an observation of the whole net.
\end{itemize}

\subsection{Policy network}\label{sec:network}
The policy and critic are graph neural networks over the observation graph. We use
relational attention with basis-decomposed relation weights, in the style of
R-GCN~\cite{rgcn2018} and RGAT~\cite{rgat2019}, with the design choices of
HGT~\cite{hu2020hgt} (per-type queries, gated skip connections, a learned prior per
relation). This is not a contribution of the paper, but it matters: with a per-type
HGT, PPO fails on our multi-site benchmark, while with this network it reaches the
heuristic (Sect.~\ref{sec:results}). Message passing follows the arcs of the net, so
condition (C2) holds.
